% Options for packages loaded elsewhere
\PassOptionsToPackage{unicode}{hyperref}
\PassOptionsToPackage{hyphens}{url}
\PassOptionsToPackage{dvipsnames,svgnames,x11names}{xcolor}
%
\documentclass[
  10pt,
]{article}
\usepackage{amsmath,amssymb}
\usepackage{iftex}
\ifPDFTeX
  \usepackage[T1]{fontenc}
  \usepackage[utf8]{inputenc}
  \usepackage{textcomp} % provide euro and other symbols
\else % if luatex or xetex
  \usepackage{unicode-math} % this also loads fontspec
  \defaultfontfeatures{Scale=MatchLowercase}
  \defaultfontfeatures[\rmfamily]{Ligatures=TeX,Scale=1}
\fi
\usepackage{lmodern}
\ifPDFTeX\else
  % xetex/luatex font selection
  \setmainfont[]{DejaVu Serif}
  \setsansfont[]{DejaVu Sans}
  \setmonofont[]{DejaVu Sans Mono}
  \setmathfont[]{Latin Modern Math}
\fi
% Use upquote if available, for straight quotes in verbatim environments
\IfFileExists{upquote.sty}{\usepackage{upquote}}{}
\IfFileExists{microtype.sty}{% use microtype if available
  \usepackage[]{microtype}
  \UseMicrotypeSet[protrusion]{basicmath} % disable protrusion for tt fonts
}{}
\makeatletter
\@ifundefined{KOMAClassName}{% if non-KOMA class
  \IfFileExists{parskip.sty}{%
    \usepackage{parskip}
  }{% else
    \setlength{\parindent}{0pt}
    \setlength{\parskip}{6pt plus 2pt minus 1pt}}
}{% if KOMA class
  \KOMAoptions{parskip=half}}
\makeatother
\usepackage{xcolor}
\usepackage[margin=2.0cm]{geometry}
\usepackage{longtable,booktabs,array}
\usepackage{calc} % for calculating minipage widths
% Correct order of tables after \paragraph or \subparagraph
\usepackage{etoolbox}
\makeatletter
\patchcmd\longtable{\par}{\if@noskipsec\mbox{}\fi\par}{}{}
\makeatother
% Allow footnotes in longtable head/foot
\IfFileExists{footnotehyper.sty}{\usepackage{footnotehyper}}{\usepackage{footnote}}
\makesavenoteenv{longtable}
\setlength{\emergencystretch}{3em} % prevent overfull lines
\providecommand{\tightlist}{%
  \setlength{\itemsep}{0pt}\setlength{\parskip}{0pt}}
\setcounter{secnumdepth}{-\maxdimen} % remove section numbering
% doc_header.tex -- LaTeX preamble additions used by build_docs.sh for every
% Markdown -> LaTeX -> PDF conversion (supplements, README, submission files).
\usepackage{etoolbox}
% longtable may not be loaded by pandoc when a document has no table
\usepackage{longtable}
% bold vectors: \bm -> \symbf under unicode-math (XeLaTeX), bm package otherwise
\makeatletter
\@ifpackageloaded{unicode-math}{\providecommand{\bm}[1]{\symbf{#1}}}{\usepackage{bm}}
\makeatother
% smaller type in long tables so that wide audit tables fit the text width;
% inside a table a line break is allowed after an underscore of plain text
% (pandoc writes it as \_), so that long identifiers wrap inside their column
% rows are set a little apart, so that symbols with indices do not touch
\AtBeginEnvironment{longtable}{\footnotesize\renewcommand{\arraystretch}{1.45}%
  \let\sptunderscore\_\def\_{\sptunderscore\allowbreak}}
\setlength{\LTcapwidth}{\textwidth}
% allow line breaks after slashes and hyphens inside table cells
\usepackage{xurl}
% code blocks: break over-long lines instead of running into the margin.
% Plain ``` blocks are typeset by pandoc as {verbatim}, highlighted ones
% (```bash) as {Highlighting}; both are re-declared as breakable fancyvrb
% environments.  The continuation line is marked with a small hooked arrow.
\usepackage{fvextra}
\RecustomVerbatimEnvironment{verbatim}{Verbatim}{breaklines,breakanywhere,fontsize=\small}
\makeatletter
\@ifundefined{Highlighting}{}{%
  \RecustomVerbatimEnvironment{Highlighting}{Verbatim}{breaklines,breakanywhere,fontsize=\small,commandchars=\\\{\}}}
\makeatother
% inline code (pandoc writes it as \texttt{...}): never hyphenate -- an inserted
% hyphen could be read as part of a file name -- but allow a line break after
% an underscore, slash, comma, colon or equals sign, so that long paths and
% column lists wrap inside the text width and inside table columns.
% Inside a formula (e.g. the name of a notebook symbol in an equation) the same
% command must first leave math mode, which \text of amsmath does.
\ExplSyntaxOn
\tl_new:N \l__spt_code_tl
\cs_new_protected:Npn \__spt_texttt:n #1
  {
    \tl_set:Nn \l__spt_code_tl {#1}
    \regex_replace_all:nnN { ([\/\,\:\=]) } { \1 \c{allowbreak} } \l__spt_code_tl
    \regex_replace_all:nnN { \c{_} } { \c{_} \c{allowbreak} } \l__spt_code_tl
    \group_begin: \ttfamily \hyphenchar \font = -1 \scan_stop: \l__spt_code_tl \group_end:
  }
\RenewDocumentCommand \texttt { m }
  {
    \mode_if_math:TF { \text { \__spt_texttt:n {#1} } } { \__spt_texttt:n {#1} }
  }
\ExplSyntaxOff
% A figure is written in the Markdown source as an image paragraph followed by
% a caption paragraph ("**Figure S1.** ..."), so that the caption is also
% visible on GitHub.  An image that starts a paragraph is centered.
\usepackage{graphicx}% pandoc loads it only for documents that contain an image
\NewCommandCopy\sptincludegraphics\includegraphics
\RenewDocumentCommand\includegraphics{O{} m}{%
  \ifvmode\begin{center}\sptincludegraphics[#1]{#2}\end{center}%
  \else\sptincludegraphics[#1]{#2}\fi}
% keep figures where they are referenced
\usepackage{float}
\floatplacement{figure}{H}
% running header with the document title
\usepackage{fancyhdr}
\pagestyle{fancy}
\fancyhf{}
\fancyhead[L]{\small PROJECT spacetimes --- documentation and provenance}
\fancyhead[R]{\small\thepage}
\renewcommand{\headrulewidth}{0.4pt}
% 16x16 matrices: tighter columns so that a full matrix fits the text width
\setlength{\arraycolsep}{3.2pt}
% The documents of this project contain long tagged symbols, file paths and
% identifiers that cannot be hyphenated.  A paragraph is therefore allowed to
% have looser lines before TeX gives up and lets a line stick out of the margin.
\AtBeginDocument{\tolerance=2000 \emergencystretch=5em \hbadness=2000}
\ifLuaTeX
  \usepackage{selnolig}  % disable illegal ligatures
\fi
\IfFileExists{bookmark.sty}{\usepackage{bookmark}}{\usepackage{hyperref}}
\IfFileExists{xurl.sty}{\usepackage{xurl}}{} % add URL line breaks if available
\urlstyle{same}
\hypersetup{
  pdftitle={P9 --- Error log of PROJECT spacetimes},
  pdfauthor={Prepared for Patrick Lee Nash, Ph.D.},
  colorlinks=true,
  linkcolor={blue},
  filecolor={Maroon},
  citecolor={blue},
  urlcolor={blue},
  pdfcreator={LaTeX via pandoc}}

\title{P9 --- Error log of PROJECT spacetimes}
\usepackage{etoolbox}
\makeatletter
\providecommand{\subtitle}[1]{% add subtitle to \maketitle
  \apptocmd{\@title}{\par {\large #1 \par}}{}{}
}
\makeatother
\subtitle{Provenance: every error made during the work, how it was
found, and how it was corrected}
\author{Prepared for Patrick Lee Nash, Ph.D.}
\date{November 12, 2026}

\begin{document}
\maketitle

\hypertarget{purpose}{%
\section{Purpose}\label{purpose}}

Hard Rule 7 of the task specification forbids shortcuts and demands that
problems be explained. This file therefore records every error that was
made while the project was carried out, including errors that were
corrected before anything was published. For each error it states what
was wrong, how it was detected, what was done, and where the corrected
statement can be checked. An error that is not listed here and is found
later is to be added, not silently repaired.

\hypertarget{errors-in-the-construction-and-verification-of-the-dirac-matrices}{%
\section{Errors in the construction and verification of the Dirac
matrices}\label{errors-in-the-construction-and-verification-of-the-dirac-matrices}}

\begin{longtable}[]{@{}
  >{\raggedright\arraybackslash}p{(\columnwidth - 8\tabcolsep) * \real{0.0435}}
  >{\raggedright\arraybackslash}p{(\columnwidth - 8\tabcolsep) * \real{0.3043}}
  >{\raggedright\arraybackslash}p{(\columnwidth - 8\tabcolsep) * \real{0.2319}}
  >{\raggedright\arraybackslash}p{(\columnwidth - 8\tabcolsep) * \real{0.2464}}
  >{\raggedright\arraybackslash}p{(\columnwidth - 8\tabcolsep) * \real{0.1739}}@{}}
\toprule\noalign{}
\begin{minipage}[b]{\linewidth}\raggedright
No.
\end{minipage} & \begin{minipage}[b]{\linewidth}\raggedright
What was wrong
\end{minipage} & \begin{minipage}[b]{\linewidth}\raggedright
How it was found
\end{minipage} & \begin{minipage}[b]{\linewidth}\raggedright
Correction
\end{minipage} & \begin{minipage}[b]{\linewidth}\raggedright
Where to check
\end{minipage} \\
\midrule\noalign{}
\endhead
\bottomrule\noalign{}
\endlastfoot
E1 & The first version of check \texttt{71-P4} stated the block form of
the complex projectors of 7+1 with the wrong sign of the off-diagonal
blocks. & The program compares the statement with the computed matrices;
the check failed. & The statement was corrected: \(\chi_{1671}\) has the
blocks \((0,-i;\,i,0)\), and \(2P_{1671}^{L}\) has the blocks
\((1,i;\,-i,1)\). & \texttt{data/matrix\_checks.txt}, line
\texttt{71-P4} \\
E2 & The 7+1 matrices of the project were first defined as the covariant
matrices \texttt{Γ17{[}swap04(A){]}} of \texttt{7+1.nb}. The Dirac
equation of the notebook, however, is written with the contravariant
matrices \texttt{Γ17up}. & Re-reading sections 5.3 and 7.1 of
\texttt{7+1.nb} while writing the Dirac equation of dirac1671. & The
project's matrices are \(\gamma_{1671}^A=\)
\texttt{Γ17up{[}swap04(A){]}}; then \(\omega_{1671}\) equals
\texttt{Γ17{[}8{]}} exactly and \(\chi_{1671}\) equals \texttt{χ17}. &
\texttt{data/matrix\_checks.txt}, lines \texttt{71-N2}, \texttt{71-V1},
\texttt{71-P1} \\
E3 & Check \texttt{71-Q9} first claimed that the two-particle content of
\(\psi^TC_{1671}^{-}\gamma_{1671}^j\psi\) is \(+\tau_{1671}^j\). & The
check failed. & The content is \(-\tau_{1671}^j\). & line
\texttt{71-Q9} \\
E4 & Check \texttt{71-P5} first claimed that the unitary matrix
\(U_{1671}^{c}\) maps \(P_{1671}^{L}\) to \(P_{1644}^{L}\). & The check
failed. & It maps \(P_{1671}^{L}\) to \(P_{1644}^{R}\) and
\(P_{1671}^{R}\) to \(P_{1644}^{L}\); the reason (two different sign
conventions for the chirality matrix in \texttt{7+1.nb}) is stated in
the check and in Sec. 2.6 of ``dirac matrices''. & line
\texttt{71-P5} \\
E5 & A commit was created although one check (E4) still failed, because
the commands were chained in a way that did not stop on failure; its
message claimed that all checks pass. & Reading the output after the
commit. & The check was corrected, the program was run again (all checks
passed), and the commit was amended before it left the workspace; it was
never pushed in the faulty state. Since then a commit is made only if
the check program returns success. & \texttt{git\ log}: commit ``Dirac
equations (dirac1644, dirac1671) and Pauli classification of fermion
pairs'' \\
E6 & The first implementation of the exact rank (fraction-free
elimination with 128-bit integers) overflowed for 256×256 matrices. &
The Rust program is compiled with overflow checks and stopped with an
error instead of returning a wrong number. & Ranks are computed modulo
the prime \(2^{61}-1\), with the argument (in the documentation of
\texttt{rank}) that a rank found modulo a prime is a lower bound of the
rank over the rational numbers, so that a full rank is proved;
statements about a null space are in addition backed by an explicit null
vector. The Python script confirms the same statements by other methods
(Gram matrices, enumeration, exact rational elimination). &
\texttt{code/spacetimes-rs/src/imat.rs};
\texttt{data/matrix\_checks\_independent.txt} \\
E7 & \textbf{Symbol collision introduced by this project.} Until
November 12, 2026 the program and the drafts used the names
\texttt{g44}, \texttt{g71}, \texttt{eta44}, \texttt{eta71},
\texttt{gamma44}, \texttt{gamma71}, \texttt{chi71} and the like for the
project's matrices. The author's notebooks use these very names for
other objects: \texttt{g44} and \texttt{g71} are metric tensors of the
gravity notebooks, \texttt{η44} is a 4×4 metric in
\texttt{Dirac\_Nash\_Identity.nb}, and \texttt{η71}, \texttt{Γ71},
\texttt{χ71}, \texttt{PL71}, \texttt{PR71} of \texttt{7+1.nb} differ
from the project's 7+1 metric, matrices, chirality matrix and
projectors. This is exactly the kind of ambiguity that Hard Rule 5
forbids. & A systematic search of all notebooks for names carrying the
tags 44, 71 and 17, made while the symbol table of ``dirac matrices''
was written. & All project symbols were renamed with the tags 1644 and
1671, which occur in no original file (source code, check statements,
data files, documents). Other letters with two meanings were separated
as well: the 8×8 matrix \(\sigma\) became \(\sigma_{88}\), the
permutation ``π'' became \(\mathrm{swap}_{04}\), the 8×8 blocks ``ρ'' of
7+1 became \(\tau_{1671}^A\), the multi-index ``K'' was replaced by
explicit index sets, and rotation parameters are written \(\theta_{AB}\)
instead of \(\omega_{AB}\). & MASTER provenance P1; Sec. 0.2 of ``dirac
matrices'' \\
E8 & In the first draft of ``dirac matrices'' the proof of Lorentz
covariance used the transformation \(\psi'(x)=S^{-1}\psi(\Lambda x)\)
and contracted the indices of \(\Lambda\) incorrectly. & Re-deriving
every formula of the draft before the document was assembled. & The
correct statement is \(\psi'(x)=S\,\psi(\Lambda x)\) with
\(S\gamma^AS^{-1}=\Lambda^A{}_B\gamma^B\); the derivation is given in
Sec. 1.8.2. The faulty draft was never committed. & ``dirac matrices'',
Sec. 1.8.2 \\
E9 & The first version of the routine of the Python script that counts
invariant pair matrices mixed the linear equations of different
generators. & Review of the code before its first complete run. & Each
generator contributes its own 64 equations for the 28 unknowns. The
corrected routine reproduces the numbers of the Rust program (0 and 10
for 7+1; 6, 10, 1 and 0 for 4+4). The faulty version was never
committed. & \texttt{code/scripts/verify\_matrices\_independent.py},
function \texttt{invariant\_dimension} \\
E10 & The draft of Sec. 1.6 said that cell 166 of \texttt{4+4.nb}
verifies all five projector relations; the cell verifies three of them.
& Comparing the sentence with the cell. & The text now says which three
relations the notebook verifies; all five are verified by checks
\texttt{44-P1}, \texttt{44-P2}. Never committed in the faulty form. &
``dirac matrices'', Sec. 1.6 \\
\end{longtable}

\end{document}
